\subsection{Population-Level Analysis Confirms Learning in both Delay and Trace Conditioning}





\subsection{Population-Level Analysis Confirms Learning in 3 s Trace Conditioning}

% Define trace conditioning intervals
In subsequent experiments, we investigated whether larvae could form associations between temporally separated events using trace conditioning protocols. Under these protocols, during Train, CS was temporally separated from US by \SI{3}{\second} (3sTrace) or \SI{10}{\second} (10sTrace; \FIG{fig1} XXXX).


% region Pop. Heatmaps
\paragraph{Population Heatmaps}

% todo 3sTrace

To capture these trends, we aligned and aggregated scaled vigor data from multiple fish for each experimental condition relative to CS onset (Materials and Methods).

the CR first appeared near CS offset, close to the expected US time. As conditioning progressed, the CR shifted earlier to occur shortly after CS onset. During testing, the CR was evident in early trials but extinguished over time, leaving only a mild, startle-like response similar to controls.


When the trace interval was extended to \SI{10}{\second}, conditioning became substantially more difficult. Heatmaps of pooled scaled vigor revealed only a subtle suppression that developed after several conditioning trials, localized primarily within the trace interval (between CS offset and US onset). This weak and delayed reduction---absent in control fish---suggests that some associative learning can occur even across this extended temporal gap, though with greatly diminished strength.
%! The 10sTrace heatmaps lacked the clear CS-locked suppression seen in shorter traces, showing instead only a modest reduction in vigor during the trace interval.

The UR remained evident during training as a bright band of elevated vigor (\textasciitilde\SI{9}{\second} after CS onset in Delay, Sup. Fig. for US-aligned data...).



% Transition from single-fish examples to population analysis
After qualitatively assessing individual learning trajectories, we next analyzed responses at the population level.
To determine whether the response patterns observed in individual fish were consistent across animals, we pooled CS-aligned tail-movement data for each conditioning group.


% !!!! vigor is undefined in the absence of movement (Materials and Methods). the heatmaps display the average vigor within each time bin.
Because vigor was undefined when a fish was not moving, each heatmap shows the mean scaled vigor of the fish that exhibited tail movements within each time bin. The heatmaps therefore describe the intensity of movement among active fish rather than the proportion of fish that moved.

Because vigor is undefined in the absence of movement (Materials and Methods), the heatmaps display the average scaled vigor across fish that exhibited tail movements within each time bin. Consequently, bins in which only a small subset of fish moved represent the average scaled vigor of that active subset, rather than population-wide immobility. Periods when most animals were immobile---such as the post-US immobility phase of the UR following the initial optovin-evoked burst---therefore do not appear uniformly dark, but instead reflect the residual activity of the few individuals that remained moving.

% Delay and 3sTrace
In the Delay and 3sTrace groups (\FIG{fig2} A,B), CS-US pairing produced a progressive reduction in vigor during training. The CR initially emerged near the expected US time and then shifted earlier during conditioning, becoming apparent soon after CS onset. This suppression persisted into the first test trials but weakened during later testing, consistent with extinction.
% todo show heatmaps with percentage of active fish per time bin (Sup. Fig. ???)


% Control
The corresponding control groups showed no sustained CS-locked suppression. Their CS-aligned vigor remained broadly stable across trials, apart from a brief increase during CS presentation. Because US delivery was explicitly unpaired and varied in time across control animals, US-evoked responses were not aligned in the pooled CS-aligned heatmaps.

The unpaired control groups provided a critical baseline for comparison.
These controls showed no sustained, learning-dependent changes.
CS-aligned scaled vigor in control animals remained relatively constant throughout the entire experiment, apart from a modest but consistent transient increase during CS presentation (visible as a brief intensity rise in the heatmaps). The UR was not apparent in pooled CS-aligned control data because each fish experienced a distinct control protocol, in which US presentations were explicitly unpaired and temporally varied. Consequently, pooling CS-aligned data across control individuals temporally misaligned the URs, causing them to appear diffuse or attenuated in the population heatmaps.
US-aligned scaled vigor remained stable in the period preceding the US, showing no systematic anticipatory modulation---activity (lacked any build-up or suppression leading up to US onset). 
% todo Sup. Fig. with US-aligned data for control
Therefore, the training phase in the control group of fish did not induce consistent activity locked to CS onset nor US onset across trials (Fig. ??? and Sup. Fig. ???).

\paragraph{Control}
Across all experimental conditions, the unpaired control groups provided a critical baseline for comparison.


% 10sTrace
The 10sTrace group showed a different pattern. Pooled heatmaps revealed, at most, a weak suppression during the trace interval between CS offset and US onset. This effect was not consistently distinguishable from control responses at the population level, suggesting that the 10-s interval substantially reduced the strength of associative learning.

%! The 10sTrace heatmaps lacked the clear CS-locked suppression seen in shorter traces, showing instead only a modest reduction in vigor during the trace interval.






\paragraph{Block-Level Statistics and Extinction}



Our quantitative analysis focused on comparing vigor change at CS onset across three distinct phases, each representing a 5-trial block:
Pre-Train (PT): The final 5 trials of pre-training, establishing baseline responsiveness.
Early Test (ET): The first 5 trials of testing, measuring the expression of the learned CR.
Late Test (LT): The final 5 trials of testing, assessing the extinction of the CR.
% todo Change abbreviations of experimental blocks of trials



% Quantify Delay-group CR expression using predefined 5-trial blocks.
Block-level analysis of the average vigor change at CS onset confirmed robust CR expression in the Delay group. During Early Test (first 5 test trials), vigor significantly decreased at CS onset relative to both the Pre-Train baseline (final 5 pre-train trials) and the Control group. This suppression was transient: by Late Test (final 5 test trials), vigor change returned toward pre-train levels, indicating extinction of the learned response.
% vigor change at CS onset (still need to find a better name for this)






\paragraph{Vigor Change at CS Onset Across Trials}
% Introduce trial-by-trial NV to track the learning trajectory and extinction dynamics.
To further characterize the learning trajectory and extinction dynamics, we quantified vigor change on a trial-by-trial basis.
% Describe trial-by-trial CR emergence, timing refinement, persistence, and extinction.
Trial-by-trial average vigor change at CS onset showed gradual CR acquisition in the Delay group during conditioning. Suppression within the CS--US interval emerged over approximately the first 10 conditioning trials (around overall trial 20). This suppression persisted into the initial testing phase (approximately trials 61--75) and declined across later test trials, consistent with extinction. In contrast, the control group remained near baseline across trials and showed no consistent trial-locked suppression pattern.

% todo Explain choice of CR windows in Materials and Methods (based on Figure 3)

%! todo when saying "rapid" learning, analyze if after a spurious CS-US pairing there is a sudden drop in NV




% This will be mostly in Sup. Fig.?

\paragraph{Vigor Change in US-aligned Data}
% US aligned Control Data
To confirm that the learned response was specifically an anticipatory behavior cued by the CS (and not other environmental stimuli), we analyzed control vigor change at CS onset aligned to the US onset during conditioning.
Control groups for both the Delay and 3sTrace conditions showed vigor change at CS onset fluctuating around 1.0, with no systematic activity patterns emerging before their randomly timed US presentations.
This behavioral stability in US-aligned control data (also exemplified in pooled scaled vigor heatmaps) strongly indicates that only the explicitly presented CS became predictive of the US. This finding ensures that the learned behavioral changes observed in conditioning groups can be confidently attributed to the specific CS-US contingencies.




\paragraph{Control Specificity: Violet Light Without Optovin}
% Specificity control: test whether violet light pulse alone can serve as an effective US.
In a preliminary control experiment, fish experienced Delay-style CS-US pairing with \SI{100}{\milli\second} violet-light pulses but without optovin. Under these conditions, violet light alone did not function as an effective US (Sup. Fig. 3), supporting the specificity of optovin-mediated reinforcement.

\paragraph{Interpretive Bridge: Baseline Stabilization and UR Profile}
% Keep concise interpretation here; detailed mechanistic interpretation can be expanded in Discussion.
The optovin-evoked UR appeared biphasic, with an early high-amplitude motor burst followed by a longer immobility phase. This qualitative profile is compatible with a fast escape-like component followed by a delayed inhibitory state and provides context for why maintaining baseline movement was important for CR quantification.

% todo Sup. Fig. XXXX with UV light only control


% todo Sup. Fig. with US-aligned data for control


% todo show heatmaps with percentage of active fish per time bin (Sup. Fig. XXXX)


% todo Change abbreviations of experimental blocks of trials
% todo ? How much was the median suppression?







%% endregion Pop. Analysis